% Journal-style academic article template for texforge.
% Text in double braces is a placeholder filled in when the project is
% created; every % comment says what to replace.
% Tectonic reads UTF-8 directly, so no input encoding package is loaded.
\documentclass[a4paper,11pt,twoside]{article}

% Document language from the {{language}} placeholder (default english).
\usepackage[{{language}}]{babel}
\usepackage[T1]{fontenc}
\usepackage{lmodern}
\usepackage{microtype}
\usepackage[a4paper,margin=2.5cm,headheight=14pt,headsep=12pt]{geometry}
\usepackage{amsmath,amssymb}
\usepackage{booktabs}
\usepackage{float}
\usepackage{graphicx}
\usepackage{fancyhdr}
\usepackage{authblk}
\usepackage[authoryear,round]{natbib}
\usepackage{xcolor}

% One restrained accent for links; everything else is monochrome.
\definecolor{articlelink}{HTML}{1F4E79}

\usepackage[
  colorlinks=true,
  linkcolor=articlelink,
  citecolor=articlelink,
  urlcolor=articlelink,
  breaklinks=true,
  pdftitle={{{title}}},
  pdfauthor={{{author}}}
]{hyperref}

% Running head: short title on the inner edge, page number on the outer edge.
% Replace \shorttitle with a shorter form if the full title does not fit.
\newcommand{\shorttitle}{{{title}}}
\pagestyle{fancy}
\fancyhf{}
\fancyhead[LE,RO]{\thepage}
\fancyhead[LO,RE]{\small\itshape\nouppercase{\shorttitle}}
\renewcommand{\headrulewidth}{0.4pt}
\fancypagestyle{plain}{%
  \fancyhf{}%
  \fancyhead[LE,RO]{\thepage}%
  \renewcommand{\headrulewidth}{0pt}%
}

% Author block: the superscript numbers come from authblk.
% Replace the affiliation and the e-mail with your own.
\title{{{title}}}
\author[1]{{{author}}}
\affil[1]{{{affiliation}}\\ \texttt{{{email}}}}
\date{\today}

\begin{document}

  \maketitle

  \begin{abstract}
    % Replace this paragraph with a short summary of the problem, the method
    % and the main result. Four to six sentences are enough.
    This article is a worked example of a single-column journal article. It
    shows one equation, one table, one diagram and one code listing, so you can
    keep the parts you need and replace the rest. Replace this paragraph with a
    summary of your own work.
  \end{abstract}

  % Replace these with the index terms of your article.
  \noindent\textbf{Keywords:} sample term, second term, third term.

  \section{Introduction}
  % State the problem, why it matters and what this article contributes.
  This article is a complete example of a single-column journal article.
  Replace this text with the context of your work: the problem you attack,
  the evidence that the problem matters, and the contribution of the article.
  The sample method follows the textbook of~\cite{smith2022}, and a short note
  by~\cite{lee2024} covers the same topic.

  \section{Materials and Methods}
  % Describe the data, the procedure and the model that you use.
  The sample method estimates a quantity from repeated measurements. Its value
  is the arithmetic mean of the observations,
  \begin{equation}
    \bar{x} = \frac{1}{n} \sum_{i=1}^{n} x_i .
    \label{eq:mean}
  \end{equation}
  % Replace the equation above with the formula at the center of your work.
  Every result in the next section is computed from this value.

  \section{Results}
  % Report the main numbers first and explain what they mean afterwards.
  Table~\ref{tab:runs} lists three illustrative runs and the value and error
  measured for each one. The best value is shown in bold.
  % Replace the values, the caption and the label with your own measurements.
  \begin{table}[H]
    \centering
    \caption{Illustrative runs, with the measured value and error}
    \label{tab:runs}
    \begin{tabular}{lcc}
      \toprule
      Run & Value & Error \\
      \midrule
      Baseline & 1.00 & 0.05 \\
      Variant A & 1.15 & 0.07 \\
      Variant B & \textbf{1.24} & 0.06 \\
      \bottomrule
    \end{tabular}
  \end{table}

  A diagram often explains a procedure faster than text. The one below is
  drawn from its source when the document is built, so the project carries no
  image file.
  % Replace the diagram with your own; keep style= and the caption.
  \begin{mermaid}[style=editorial, caption={Sample workflow from data to result}]
    flowchart LR
    data[Collect] --> clean[Clean] --> analyze[Analyze] --> report[Report]
  \end{mermaid}

  \section{Discussion}
  % Say what the numbers mean and how they relate to earlier work.
  Say what the numbers mean, how they compare with the method
  of~\cite{doe2023}, and which limitation of the sample matters most. Keep the
  discussion tied to the result instead of repeating it.

  \section{Conclusion}
  % Restate the main result, name the limitation and the next step.
  Close with one sentence on the result, one on the limitation and one on the
  next step. Replace this sample closing with your own.

  \section{Data and Code Availability}
  % Show the few commands a reader needs to reproduce the result.
  The commands below reproduce the sample result. Replace them with the code
  or the data location of your own work.
  \begin{code}[lang=bash]
# Regenerate the sample results.
make data
make results
  \end{code}

  \bibliographystyle{plainnat}
  \bibliography{bib/references}

\end{document}
